\section{不动点与 Newton 法}
\label{sec:08-Numerical-Analysis/02-Root-Finding/02}

上一节的二分法每步只买到一个二进制位，约 \(0.30\) 个十进制位。想要十六位有效数字，得跑五十多次函数。可如果 \(f\) 光滑、导数也拿得到，那些被二分法扔掉的信息——函数值的大小、图像的走向——本来是可以换成速度的。

换法不止一种，效果差得很远。同一个方程改写成不同的迭代格式，有的三步撞上机器精度，有的原地摆动，有的很快溢出。所以先别急着写 Newton 公式。更基本的问题是：给定一条迭代规则，什么决定它收不收敛。

\subsection{把求根改写成反复代入}

求 \(f(x)=0\) 与找一个满足 \(x=g(x)\) 的点，只要 \(g\) 的构造让两边解集重合，就是同一件事。这样的 \(x\) 叫 \(g\) 的不动点，而迭代 \(x_{k+1}=g(x_k)\) 是所有迭代算法里最朴素的形式：把上一次的输出重新喂回去。

同一个方程可以有很多种改写。取 \(f(x)=x^2-2\)，至少有三种：

\[
x=\frac{2}{x},\qquad x=x^2+x-2,\qquad x=\frac12\Bigl(x+\frac{2}{x}\Bigr).
\]

三者的不动点都是 \(\pm\sqrt2\)，代数上完全等价。可从 \(x_0=1.5\) 出发跑一跑，第一种在 \(1.5\) 与 \(1.3333\) 之间来回摆动，永远停不下来；第二种走到 \(1.75,\,2.8125,\,8.72,\,82.8\)，十来步后溢出；第三种给出 \(1.416667,\,1.4142157,\,1.41421356\)，三步就把 \(\sqrt2\) 算到机器精度。

代数等价不保证数值等价。差别不在方程里，在迭代规则里。

\subsection{一个数决定收敛：\texorpdfstring{\(\abs{g'(r)}<1\)}{|g'(r)|<1}}

设 \(r=g(r)\)，记误差 \(e_k=x_k-r\)。中值定理给出

\[
e_{k+1}=g(x_k)-g(r)=g'(\xi_k)\,e_k,
\]

其中 \(\xi_k\) 夹在 \(x_k\) 与 \(r\) 之间。迭代进入根附近后 \(\xi_k\to r\)，上式趋于 \(e_{k+1}\approx g'(r)e_k\)：每步误差乘上同一个因子。因子的绝对值小于 \(1\) 就收缩，大于 \(1\) 就放大；等于 \(1\) 时一阶信息用尽，行为由更高阶的项决定。

判据只用到一个数，而且可以画出来。

\begin{center}
\begin{tikzpicture}[x=1cm,y=1cm,font=\small,>=stealth]
  \draw[->,black!55] (0,0) -- (3.7,0) node[right,font=\footnotesize] {\(x\)};
  \draw[->,black!55] (0,0) -- (0,3.7);
  \draw[black!35] (0,0) -- (3.4,3.4);
  \node[right,font=\footnotesize,black!45] at (3.05,3.45) {\(y=x\)};
  \draw[blue!75,thick] (0,3.2) -- (3.4,1.16);
  \node[below,font=\footnotesize,blue!70] at (3.0,1.28) {\(y=g(x)\)};
  \draw[red!65] (0.5,0.5) -- (0.5,2.9) -- (2.9,2.9) -- (2.9,1.46) -- (1.46,1.46)
    -- (1.46,2.324) -- (2.324,2.324) -- (2.324,1.806) -- (1.806,1.806)
    -- (1.806,2.116) -- (2.116,2.116);
  \fill[black!75] (2,2) circle (1.5pt);
  \node[right,font=\footnotesize] at (2.1,1.84) {\(r\)};
  \node[below,font=\footnotesize] at (0.5,-0.05) {\(x_0\)};
  \node[font=\footnotesize,align=center,black!70] at (1.8,-0.95) {\(\abs{g'}=0.6\)：\\螺旋向内收缩};
\begin{scope}[shift={(5.0,0)}]
  \draw[->,black!55] (0,0) -- (3.7,0) node[right,font=\footnotesize] {\(x\)};
  \draw[->,black!55] (0,0) -- (0,3.7);
  \draw[black!35] (0,0) -- (3.4,3.4);
  \node[right,font=\footnotesize,black!45] at (3.05,3.45) {\(y=x\)};
  \draw[blue!75,thick] (0.68,0.02) -- (3.05,3.575);
  \node[right,font=\footnotesize,blue!70] at (2.35,2.45) {\(y=g(x)\)};
  \draw[red!65] (1.85,1.85) -- (1.85,1.775) -- (1.775,1.775) -- (1.775,1.6625)
    -- (1.6625,1.6625) -- (1.6625,1.4938) -- (1.4938,1.4938) -- (1.4938,1.2406)
    -- (1.2406,1.2406) -- (1.2406,0.8609) -- (0.8609,0.8609) -- (0.8609,0.2914);
  \draw[->,red!65] (0.8609,0.2914) -- (0.30,0.2914);
  \fill[black!75] (2,2) circle (1.5pt);
  \node[right,font=\footnotesize] at (2.08,1.9) {\(r\)};
  \node[below,font=\footnotesize] at (1.85,-0.05) {\(x_0\)};
  \node[font=\footnotesize,align=center,black!70] at (1.8,-0.95) {\(\abs{g'}=1.5\)：\\台阶越迈越远};
\end{scope}
\end{tikzpicture}

\emph{蛛网图的读法：从 \(x_0\) 竖直走到曲线 \(y=g(x)\)，读出 \(g(x_0)\)；再水平走到对角线 \(y=x\)，把它变成下一个输入。左图斜率的绝对值小于 \(1\)，折线一圈圈缩向不动点；右图斜率大于 \(1\)，同样的走法把台阶越迈越远，起点明明就在不动点旁边。收敛条件是可以看见的：它就是这两幅折线的形状差别。}
\end{center}

回到那三种改写，在正根处算 \(g'\)：\(2/x\) 给出 \(-1\)，恰好卡在边界上，于是原地摆动；\(x^2+x-2\) 给出 \(2\sqrt2+1>1\)，注定发散；\(\frac12(x+2/x)\) 给出 \(0\)，比任何线性收缩都快。第三个正是下面要讲的 Newton 法在这个方程上的样子。

\(\abs{g'(r)}<1\) 值得多看一眼，因为它在这本书里已经出现过几次，只是换了名字。线性递推 \(u_{k+1}=Au_k\) 的收敛与否由特征根的模是否小于 \(1\) 决定，见\hyperref[sec:04-Discrete-Mathematics/02-Recurrences/02]{``线性递推与特征根''}；Markov 链趋于平稳分布的速度由转移矩阵第二大特征值的模控制，见\hyperref[sec:06-Random-Processes/03-Markov-Chains/04]{``平稳分布、可逆性与遍历收敛''}；离散动力系统的不动点是吸引还是排斥，判据仍是 \(\abs{g'}\) 与 \(1\) 的大小关系，见\hyperref[sec:15-Differential-Equations-and-Dynamical-Systems/02-Stability-and-Dynamical-Systems/03]{``离散动力系统、分岔与混沌''}。它们是同一条定理换了几套衣服：一次迭代把误差线性放大 \(\lambda\) 倍，\(n\) 次之后是 \(\lambda^n\)，命运只取决于 \(\abs\lambda\) 站在 \(1\) 的哪一侧。求根算法的收敛性分析与动力系统的稳定性分析，从来不是两门学问。

\subsection{局部判据要升级才有全局保证}

\(\abs{g'(r)}<1\) 是局部结论：它保证初值``足够近''时收敛，却说不出多近算近——而 \(r\) 本来就是你要找的东西，用它来判断初值好不好，是循环论证。要把``足够近''换成可检验的区间，条件得加强。

\begin{theorem}
设闭区间 \(I\) 满足 \(g(I)\subseteq I\)，且存在 \(q<1\) 使 \(\abs{g'(x)}\le q\) 对一切 \(x\in I\) 成立。则 \(g\) 在 \(I\) 内有唯一不动点，且从 \(I\) 内任意初值出发的迭代都收敛于它。
\end{theorem}

证明骨架只有两句：由中值定理，\(\abs{g'}\le q\) 给出 \(\abs{g(x)-g(y)}\le q\abs{x-y}\)，相邻两步的间距按几何级数衰减，序列是 Cauchy 列因而收敛；不变性 \(g(I)\subseteq I\) 保证极限没跑出 \(I\)，收缩性保证不动点唯一。这是 Banach 不动点定理的一维形态。

两个条件各防一种失效。丢掉不变性，迭代第一步就可能跳出 \(I\)，区间内的斜率控制立刻作废。丢掉一致的 \(q<1\)，只在一点成立的收缩可以在稍远处失灵：\(g(x)=x+x^3\) 在 \(0\) 处 \(g'=1\)，任何包含原点的区间上都找不到统一的 \(q<1\)，从任何非零初值出发的迭代都越走越远。

收缩性还附赠一个后验误差界：

\[
\abs{x_k-r}\le\frac{q}{1-q}\,\abs{x_k-x_{k-1}}.
\]

真根不必知道，看最近两步的间距就够了。留意前面那个系数——\(q\) 接近 \(1\) 时它会炸开。收敛很慢的时候，小步长根本不能说明离根近，这是自动停止条件里最常见的坑。

\subsection{Newton 法：用切线代替曲线}

既然 \(\abs{g'(r)}\) 越小越快，不妨直接把它做成零。

想法本身很朴素：在 \(x_k\) 处曲线不好解，就换成它的切线，解切线的零点。这是一阶 Taylor 近似的直接使用，见\hyperref[sec:02-Calculus/03-Differentiation/05]{``线性近似与极值问题''}。写下

\[
f(x_k+s)\approx f(x_k)+f'(x_k)s,
\]

令右边为零解出 \(s\)，得到

\[
x_{k+1}=x_k-\frac{f(x_k)}{f'(x_k)}.
\]

对应的迭代函数是 \(g(x)=x-f(x)/f'(x)\)。求导并在简单根处取值：

\[
g'(x)=\frac{f(x)f''(x)}{f'(x)^2},\qquad g'(r)=0\quad(f(r)=0,\ f'(r)\ne0).
\]

线性因子被构造性地消掉了。上一小节的 \(e_{k+1}\approx g'(r)e_k\) 于是退化成 \(0\)，误差的衰减只能由更高阶的项接手——二次收敛正是从这里来的。

\begin{center}
\begin{tikzpicture}[x=2.05cm,y=0.40cm,font=\small,>=stealth]
  \draw[->,black!55] (0.85,0) -- (3.35,0) node[right,font=\footnotesize] {\(x\)};
  \draw[blue!75,thick] plot[domain=1.0:3.12,samples=70] (\x,{\x*\x-2});
  \node[left,font=\footnotesize,blue!70] at (2.52,5.6) {\(y=f(x)\)};
  \draw[red!70] (1.8333,0) -- (3.12,7.72);
  \draw[red!70] (1.4621,0) -- (2.35,3.256);
  \draw[red!70] (1.4150,0) -- (1.95,1.565);
  \draw[black!30,dashed] (3,0) -- (3,7);
  \draw[black!30,dashed] (1.8333,0) -- (1.8333,1.3611);
  \fill[black!75] (3,7) circle (1.4pt);
  \fill[black!75] (1.8333,1.3611) circle (1.4pt);
  \fill[black!75] (1.4621,0.1378) circle (1.4pt);
  \fill[red!70] (3,0) circle (1.3pt);
  \fill[red!70] (1.8333,0) circle (1.3pt);
  \fill[red!70] (1.4621,0) circle (1.3pt);
  \draw[black!45,dashed] (1.4142,0.15) -- (1.4142,4.2);
  \node[above,font=\footnotesize,black!60] at (1.4142,4.2) {根 \(r=\sqrt2\)};
  \node[below,font=\footnotesize] at (3,-0.15) {\(x_0\)};
  \node[below,font=\footnotesize] at (1.8333,-0.15) {\(x_1\)};
  \node[below,font=\footnotesize] at (1.4621,-0.15) {\(x_2\)};
\end{tikzpicture}

\emph{曲线是 \(f(x)=x^2-2\)，从 \(x_0=3\) 出发。每一条红线是当前点处的切线，它与横轴的交点就是下一个猜测：\(3\to1.833\to1.462\)。第四步给出 \(1.41500\)，与真根 \(1.41421\) 的差已经小到画不出来，所以图上没有 \(x_3\)——这正是``位数翻倍''在纸面上的样子。}
\end{center}

把 \(f(x)=x^2-a\) 代进去，更新式化简成 \(x_{k+1}=\frac12(x_k+a/x_k)\)，正是前面第三种改写，也是古人开平方的办法：估计偏大时 \(a/x_k\) 偏小，取平均把误差从两侧夹回来。

\subsection{位数翻倍来自两行 Taylor}

骨架很短。设 \(r\) 是简单根，在 \(x_k\) 处把 \(f(r)\) 展到二阶并带 Lagrange 余项：

\[
0=f(r)=f(x_k)-f'(x_k)e_k+\tfrac12f''(\xi_k)e_k^2 .
\]

解出 \(f(x_k)\)，代进 \(e_{k+1}=e_k-f(x_k)/f'(x_k)\)，一阶项恰好抵消，剩下

\[
e_{k+1}=\frac{f''(\xi_k)}{2f'(x_k)}\,e_k^2 .
\]

条件对账做三笔。要 \(f\) 二阶连续可导，余项才有意义；要 \(f'(x_k)\) 不接近零，分母才不炸；要 \(x_k\) 已经落在根的某个邻域内，\(\xi_k\) 才跟着被夹住，\(f''/(2f')\) 才有统一的上界 \(C\)。三条齐备，才有 \(\abs{e_{k+1}}\le C\abs{e_k}^2\)。

这个不等式的读法比它的推导重要。误差 \(10^{-4}\) 的下一步大约是 \(10^{-8}\)，再下一步 \(10^{-16}\)：正确的十进制位数每步翻倍，三四步内撞上机器精度。所以 Newton 法的迭代次数几乎不随精度要求增长，要八位和要十六位只差一步——这与二分法每多一位就多三步半，是两种完全不同的成本曲线。

同一个不等式也写着它的软肋。\(C\abs{e_k}^2<\abs{e_k}\) 要求 \(\abs{e_k}<1/C\)。初值不在这道门槛之内时，平方项非但不帮忙，反而把误差放大。

\subsection{离根远时，切线可能指向任何地方}

反例比警告有用。取 \(f(x)=\arctan x\)，唯一的根是 \(0\)，函数处处光滑，导数处处为正，条件看起来无可挑剔。Newton 更新是

\[
x_{k+1}=x_k-(1+x_k^2)\arctan x_k .
\]

从 \(x_0=0.5\) 出发，两步就到 \(3\times10^{-4}\) 量级，干脆利落。把初值挪到 \(x_0=1.5\)，同一条公式给出 \(x_1\approx-1.694\)，\(x_2\approx2.322\)，\(x_3\approx-5.118\)，越震越远，再跑几步就溢出。分水岭在 \(\abs{x_0}\approx1.3917\)——它是方程 \(2x=(1+x^2)\arctan x\) 的正根。跨过这条线，同一个算法、同一个函数，行为从收敛翻成发散。

\begin{center}
\begin{tikzpicture}[x=1.15cm,y=1.45cm,font=\small,>=stealth]
  \draw[->,black!55] (-4.5,0) -- (4.6,0) node[right,font=\footnotesize] {\(x\)};
  \draw[->,black!55] (0,-1.8) -- (0,1.8);
  \draw[black!25,dashed] (-4.4,1.5708) -- (4.3,1.5708);
  \draw[black!25,dashed] (-4.4,-1.5708) -- (4.3,-1.5708);
  \node[right,font=\footnotesize,black!55] at (3.2,1.68) {\(\pi/2\)};
  \draw[blue!75,thick] plot[domain=-4.3:4.3,samples=90] (\x,{rad(atan(\x))});
  \draw[red!70] (-1.694,0) -- (2.3,1.229);
  \draw[red!70] (2.322,0) -- (-2.6,-1.272);
  \draw[->,red!70] (2.322,1.1642) -- (-4.3,0.128);
  \draw[black!30,dashed] (1.5,0) -- (1.5,0.9828);
  \draw[black!30,dashed] (-1.694,0) -- (-1.694,-1.0378);
  \draw[black!30,dashed] (2.322,0) -- (2.322,1.1642);
  \fill[black!75] (1.5,0.9828) circle (1.4pt);
  \fill[black!75] (-1.694,-1.0378) circle (1.4pt);
  \fill[black!75] (2.322,1.1642) circle (1.4pt);
  \fill[red!70] (1.5,0) circle (1.3pt);
  \fill[red!70] (-1.694,0) circle (1.3pt);
  \fill[red!70] (2.322,0) circle (1.3pt);
  \node[below,font=\footnotesize] at (1.42,-0.06) {\(x_0\)};
  \node[below,font=\footnotesize] at (-1.98,-0.06) {\(x_1\)};
  \node[below,font=\footnotesize] at (2.45,-0.06) {\(x_2\)};
  \node[right,font=\footnotesize,black!70] at (-4.3,0.74) {下一步落在 \(x_3\approx-5.1\)，已经出图};
\end{tikzpicture}

\emph{与上一张图成对来看。同样是``画切线、取交点''，这里每一步都跨过根落到对面更远处：\(1.5\to-1.694\to2.322\to-5.118\)。原因写在图形上——\(\arctan\) 在远处几乎水平，切线的斜率趋于零而函数值不趋于零，步长 \(-f/f'\) 于是大得离谱。快和失控是同一个机制的两面。}
\end{center}

导数接近零是最常见的触发方式，此外还有：切线指向定义域之外，函数在那里根本没有值；或者迭代在若干点之间形成闭环，既不收敛也不发散。共同点在于 Newton 法不维护任何不变量。二分法每步之后``根在当前区间内''依然成立，随时可以出示；Newton 法迭代到一半，你手里没有任何东西能证明根还在附近。

速度与保证在这里正面对撞。二分法慢，但只要 \(f\) 连续、端点异号，它从不失手，所需步数还能在开跑前算出来。Newton 法快得多，而换来的这份快要靠三条额外假设撑着——可导、导数不接近零、初值足够近——最后一条恰恰最难在事前检验。

\subsection{重根把切线压平，二次收敛塌回线性}

还有一种失效不是发散，而是悄悄变慢。设 \(f(x)=(x-r)^mh(x)\)，\(h(r)\ne0\)，重数 \(m>1\)。根附近 \(f/f'\approx(x-r)/m\)，于是

\[
g'(r)=1-\frac1m=\frac{m-1}{m}\ne0 .
\]

线性因子回来了。二重根每步误差减半，跟二分法一个速度；三重根每步只减掉三分之一，比二分法还慢。

症状可以直接观察：打印相邻两步的误差比 \(\abs{e_{k+1}}/\abs{e_k}\)。二次收敛时这个比值飞快趋于零，线性收敛时它稳定在 \(0.5\) 或 \(2/3\) 附近纹丝不动。看到一个稳住不降的比值，基本可以断定踩上了重根，或者 \(f'\) 在根处退化。重数已知时把步长乘以 \(m\) 就能恢复二次收敛，可惜重数通常不知道。

几何解释是切线被压平了。重根处 \(f'(r)=0\)，根附近导数很小，切线几乎与横轴平行，交点位置对函数值的微小变化极其敏感。

减速并非算法的缺陷。重根本身就是病态的：给 \(f\) 加一个 \(\varepsilon\) 的扰动，二重根会分裂成相距 \(\sqrt\varepsilon\) 的两个根，扰动被开方地放大——\(\varepsilon=10^{-16}\) 换来 \(10^{-8}\) 的根位移。求根问题在这一点的条件数已经坏了，见\hyperref[sec:08-Numerical-Analysis/01-Floating-Point-and-Error/03]{``条件数衡量问题本身的敏感性''}。算法慢下来只是这件事的回声：任何方法在这里都拿不到十六位有效数字，因为那些位在数据里本来就不存在。

\subsection{给快方法配护具}

纯 Newton 公式是引擎，不是车。可靠的求根器都把它装在一个框架里。

最常见的框架是借二分法的括区当安全网：维护一个异号区间，每步先算 Newton 候选点，若它落在区间内且能让区间明显缩小就采纳，否则退回一次对折。最坏情况由对折兜底，寻常情况由 Newton 推进。另一种是阻尼，把步长乘以 \(\alpha_k\in(0,1]\)，从 \(1\) 起试，不满足下降条件就减半。两者的共同点是不再无条件相信切线，先要它交出一点证据。

多元情形下切线换成切平面，除法换成解线性方程组：

\[
J_f(x_k)\,s_k=-f(x_k),\qquad x_{k+1}=x_k+s_k,
\]

其中 \(J_f\) 是 Jacobian 矩阵。实现上要解这个方程组，而不是先求出 \(J_f^{-1}\) 再乘：显式求逆运算量更大，数值稳定性也更差。局部二次收敛依然成立，门槛却更高了——\(J_f\) 在根处非奇异，初值落在一个通常很小的吸引域内。维数一高，``初值足够近''几乎无法保证。信赖域的做法是先划定一个半径，只在线性模型可信的范围内取步，再按实际下降与预测下降的比值动态调整这个半径。

这些都不只是求根的事。无约束优化的一阶最优性条件 \(\nabla f(x)=0\) 本身就是一个 \(n\) 元方程组的求根问题；对它用 Newton 法，Jacobian 就是 Hessian，得到的正是二阶优化方法，见\hyperref[sec:09-Convex-Optimization/06-Optimization-Algorithms/02]{``Newton 法用曲率校正方向''}。深度学习里很少直接使用它，理由和本节说的完全一致：高维非凸地形上初值离驻点很远，Hessian 可能不定甚至奇异，纯 Newton 步会朝鞍点甚至更高处走去，参见\hyperref[sec:09-Convex-Optimization/08-Nonconvex-and-Stochastic-Optimization/01]{``下降引理在非凸地形中只保证接近驻点''}。所有加在 Newton 法上的护具——阻尼、线搜索、信赖域、给 Hessian 加对角项——都在做同一件事：承认二次模型只在局部可信，然后为``局部''划一条线。

到这里，Newton 法的取舍已经清楚：一次导数换二次收敛，前提是可导、导数不退化、初值够近，失手时没有安全网。可导数常常拿不到——函数是一段仿真代码，或者一次远程调用，写出解析导数的成本和重写整个模型差不多。下一节把导数换成最近两次函数值的差商，看看这笔更便宜的交易能保住多少速度，又会新添哪些风险。
